\mxunit{automatismes-ch10}{Géométrie repérée~: synthèse vecteurs / droites}

\begin{fluencysheet}{mar.~30~mars~2027}{Vecteurs et colinéarité~: réinvestissement}
\mxfq{Échauffement · Flash}{GEO}{$\vec u\vc{3}{-2}$ et $\vec v\vc{-1}{5}$. Calculer $\vec u+\vec v$ et $2\vec u$.}{$\vc{2}{3}$\,; $\vc{6}{-4}$.}{}
\mxfq{Cœur · Méthode}{GEO}{Les points $A(-1;2)$, $B(1;3)$ et $C(5;5)$ sont-ils alignés\,?}{$\vect{AB}\vc21$, $\vect{AC}\vc63$~: $2\times3-1\times6=0$~: alignés.}{}
\mxfq{Cœur + · Méthode}{GEO}{$A(-1;2)$, $B(3;1)$, $C(4;4)$. Déterminer $D$ tel que $ABCD$ soit un parallélogramme.}{$\vect{AB}=\vect{DC}$~: $D=A+\vect{BC}=(-1+1;2+3)=(0;5)$.}{}
\mxfq{Retour · Flash}{PRO}{Dans une classe de $24$ élèves, $18$ sont demi-pensionnaires. Quelle proportion en pourcentage\,?}{$\dfrac{18}{24}=0{,}75=75\,\%$.}{}
\mxfq{Fil rouge · Prouver}{LIT}{Montrer que $(n+1)^2-(n-1)^2=4n$ pour tout entier $n$. Que peut-on en conclure pour la différence des carrés de deux impairs consécutifs ($n$ pair)\,?}{$(n^2+2n+1)-(n^2-2n+1)=4n$. Pour $n=2k$ pair, la différence de deux carrés d'impairs consécutifs vaut $8k$~: multiple de $8$.}{Les deux entiers $n-1$ et $n+1$ sont impairs lorsque $n$ est pair.}
\end{fluencysheet}

\begin{fluencysheet}{jeu.~1\textsuperscript{er}~avr.~2027}{Déterminant et équations de droites~: réinvestissement}
\mxfq{Échauffement · Flash}{GEO}{Pente de la droite $3x+2y-6=0$\,?}{$y=-\dfrac32x+3$~: pente $-\dfrac32$.}{}
\mxfq{Cœur · Méthode}{GEO}{Équation de la droite $(AB)$ avec $A(1;2)$ et $B(4;8)$.}{Pente $\dfrac{8-2}{4-1}=2$\,; $y=2x+p$, $2=2+p$~: $y=2x$\,; ou $2x-y=0$.}{}
\mxfq{Cœur + · Méthode}{GEO}{Les droites $2x-y+1=0$ et $4x-2y+5=0$ sont-elles parallèles, confondues ou sécantes\,?}{Vecteurs directeurs $\vc12$ et $\vc24$ colinéaires~: parallèles\,; $(0;1)$ vérifie la première mais $4\times0-2+5=3\neq0$~: strictement parallèles.}{}
\mxfq{Retour · Flash}{PUI}{Calculer $10^{-2}\times5\times10^4$.}{$5\times10^{2}=500$.}{}
\mxfq{Fil rouge · Méthode}{EQ}{Résoudre dans $\R$~: $7-2(x-1)=3x-6$.}{$x=3$}{Piège~: $-2(x-1)=-2x-2$ (erreur de signe).}
\end{fluencysheet}

\begin{fluencysheet}{lun.~5~avr.~2027}{Problèmes de géométrie plane~: choisir une méthode}
\mxfq{Échauffement · Flash}{GEO}{Distance $AB$ pour $A(0;0)$ et $B(5;12)$\,?}{$\sqrt{25+144}=13$.}{}
\mxfq{Cœur · Méthode}{DEM}{$A(0;0)$, $B(4;2)$, $C(-2;4)$. Nature du triangle $ABC$\,?}{$AB^2=20$, $AC^2=20$, $BC^2=36+4=40=AB^2+AC^2$~: rectangle isocèle en $A$.}{}
\mxfq{Cœur + · Méthode}{DEM}{Citer deux méthodes pour démontrer que $ABCD$ est un parallélogramme à l'aide de coordonnées.}{(i) $\vect{AB}=\vect{DC}$\,; (ii) les diagonales $[AC]$ et $[BD]$ ont le même milieu.}{}
\mxfq{Retour · Lecture}{LEC}{Dans une classe de $30$ élèves, un diagramme en bâtons indique $x$ élèves pour la note $12$ et $2x$ élèves pour la note $14$\,; les autres notes regroupent $18$ élèves. Trouver $x$.}{$x+2x+18=30$, donc $3x=12$ et $x=4$.}{}
\mxfq{Fil rouge · Flash}{SOL}{Une pyramide a pour base un carré de côté $6$ cm et pour hauteur $5$ cm. Calculer son volume.}{$V=\dfrac13\times6^2\times5=60$ cm$^3$.}{}
\end{fluencysheet}

\begin{fluencysheet}{mar.~6~avr.~2027}{Situations issues d'autres disciplines}
\mxfq{Échauffement · Flash}{MOD}{Un mobile a pour position $x(t)=3t+2$ (en m, $t$ en s). Où est-il à $t=4$\,?}{$x(4)=14$ m.}{}
\mxfq{Cœur · Méthode}{MOD}{Un second mobile suit $x_2(t)=5t-4$. À quel instant et où les mobiles se rencontrent-ils\,?}{$3t+2=5t-4\Leftrightarrow t=3$ s\,; position $x=11$ m.}{}
\mxfq{Cœur + · Flash}{GEO}{Un vecteur vitesse a pour coordonnées $\vc34$ (en m/s). Quelle est sa norme\,?}{$\sqrt{3^2+4^2}=5$ m/s.}{}
\mxfq{Retour · Flash}{CAL}{Calculer $\dfrac12-\dfrac35$.}{$\dfrac{5}{10}-\dfrac{6}{10}=-\dfrac1{10}$.}{}
\mxfq{Fil rouge · Rédiger}{DEM}{$ABC$ est isocèle en $A$ et $H$ est le pied de la hauteur issue de $A$. Que peut-on affirmer de $H$\,?}{$H$ est le milieu de $[BC]$ (la hauteur issue du sommet est aussi médiatrice de la base).}{}
\end{fluencysheet}

\begin{fluencysheet}{jeu.~8~avr.~2027}{Communiquer une démonstration~: implication, réciproque, équivalence}
\mxfq{Échauffement · Flash}{DEM}{Énoncer la réciproque de~: «~si un quadrilatère est un losange, alors ses diagonales sont perpendiculaires~».}{«~Si un quadrilatère a ses diagonales perpendiculaires, alors c'est un losange~».}{}
\mxfq{Cœur · Vrai/faux}{DEM}{Vrai ou faux\,? Justifier. La réciproque précédente est vraie.}{Faux. Contre-exemple~: un cerf-volant (non losange) a des diagonales perpendiculaires.}{}
\mxfq{Cœur + · Rédiger}{DEM}{L'équivalence «~$x^2=4\Leftrightarrow x=2$~» est-elle correcte\,? Corriger.}{Non~: $x=-2$ convient aussi. On a l'implication $x=2\Rightarrow x^2=4$ et l'équivalence $x^2=4\Leftrightarrow(x=2\text{ ou }x=-2)$.}{}
\mxfq{Retour · Flash}{CMP}{Comparer $2^{10}$ et $10^3$.}{$2^{10}=1\,024>1\,000=10^3$.}{}
\mxfq{Fil rouge · Flash}{REP}{Entre quels entiers consécutifs se situe $\dfrac{7}{4}$ sur une droite graduée\,? Donner sa valeur décimale.}{Entre $1$ et $2$\,; $1{,}75$.}{}
\end{fluencysheet}

\begin{fluencysheet}{lun.~26~avr.~2027}{Oral de mise en commun~: comparaison de méthodes}
\mxfq{Échauffement · Flash}{GEO}{Distance entre $A(-2;1)$ et $B(1;5)$\,?}{$\sqrt{3^2+4^2}=5$.}{}
\mxfq{Cœur · Méthode}{GEO}{Citer trois méthodes pour prouver que $A$, $B$, $C$ sont alignés.}{(i) $\deter{\vect{AB}}{\vect{AC}}=0$\,; (ii) équation de $(AB)$ vérifiée par $C$\,; (iii) $AB+BC=AC$~: égalité de distances, équivalente à $B\in[AC]$.}{}
\mxfq{Cœur + · Réflexion}{GEO}{Pour prouver que deux droites données par $y=m_1x+p_1$ et $y=m_2x+p_2$ sont parallèles, quelle méthode est la plus rapide\,?}{Comparer les pentes~: $m_1=m_2$\,; parallèles strictement si de plus $p_1\neq p_2$.}{}
\mxfq{Retour · Factoriser}{DEV}{Factoriser $6x+9$.}{$3(2x+3)$.}{}
\mxfq{Fil rouge · Méthode}{DIV}{Décomposer $360$ en facteurs premiers. Combien de diviseurs a-t-il\,?}{$360=2^3\times3^2\times5$\,; $(3+1)(2+1)(1+1)=24$ diviseurs.}{}
\end{fluencysheet}
